\chapter{Requirements Analysis}

\section{Functional Requirements}

Requirements are grouped by module and numbered for traceability. The final column
names the evidence that the requirement is met; ``test'' refers to the automated
suite described in Chapter~\ref{ch:testing}, and ``walkthrough'' to the end-to-end
demonstration in Section~\ref{sec:walkthrough}.

\begin{table}[htbp]
  \centering
  \small
  \begin{tabularx}{\textwidth}{@{}lXl@{}}
    \toprule
    \textbf{ID} & \textbf{Requirement} & \textbf{Verified by} \\
    \midrule
    \multicolumn{3}{@{}l}{\textit{Identity}} \\
    FR1 & Register with name, e-mail and password. The password is stored only as a
      BCrypt hash; the role is forced to \code{USER}. & test \\
    FR2 & Log in with e-mail and password and receive a JWT that identifies the
      account. & test, Postman \\
    FR3 & Exchange a stored JWT for the caller's profile so a relaunch restores
      the session. & test, walkthrough \\
    \multicolumn{3}{@{}l}{\textit{Vehicles and facilities}} \\
    FR4 & Add, list, update and delete own vehicles, each with a number plate and a
      class (\code{MC}, \code{LMV}, \code{HMV}). & test \\
    FR5 & Look up a facility by its code, including per-second rates for all three
      vehicle classes. & test \\
    FR6 & Create facilities, change pricing and change status through an
      administrator key. & test, walkthrough \\
    \multicolumn{3}{@{}l}{\textit{Sessions and billing}} \\
    FR7 & Start a session for one of the caller's vehicles at an active facility;
      the session begins in \code{CREATED} with no timestamps. & test \\
    FR8 & Refuse a new session when the user already holds one, or when the vehicle
      already holds one. & test \\
    FR9 & Read the live fee of a parked session, computed from the backend clock. & test \\
    FR10 & Pay for a parked session; the amount is frozen and the session moves to
      \code{PAYMENT\_PENDING}. & test, walkthrough \\
    FR11 & Cancel a session that has not entered the gate. & test \\
    FR12 & Abort \code{CREATED} sessions that are never validated at a gate. & test \\
    FR13 & Page through own session history, newest first, optionally filtered by
      status. & test \\
    \multicolumn{3}{@{}l}{\textit{Gate and devices}} \\
    FR14 & Register gate hardware; return the device secret exactly once and store
      only its hash. & test \\
    FR15 & Answer an entry verdict: open only for a \code{CREATED} session at the
      same facility, and stamp the entry time. & test \\
    FR16 & Answer an exit verdict: open only for a paid session, and allow one grace
      re-open after exit. & test \\
    FR17 & Report every verdict with a machine-readable reason, and deny with a
      reason rather than an error when the session is unknown. & test \\
    FR18 & Queue \code{OPEN} or \code{KEEP\_CLOSED} to a gate; the next poll
      consumes the command exactly once. & test \\
    \multicolumn{3}{@{}l}{\textit{Application}} \\
    FR19 & Discover a gate by BLE advertisement, list nearby facilities and fetch
      their details. & walkthrough \\
    FR20 & Send entry and exit commands over BLE with the session identifier and
      display the gate's answer, including denial reasons. & walkthrough \\
    \bottomrule
  \end{tabularx}
  \caption{Functional requirements and their verification evidence.}
  \label{tab:functional}
\end{table}

\section{Non-Functional Requirements}

\begin{table}[htbp]
  \centering
  \small
  \begin{tabularx}{\textwidth}{@{}lX@{}}
    \toprule
    \textbf{Attribute} & \textbf{Requirement and how it is satisfied} \\
    \midrule
    Security & Three credentials with separate scopes: JWT for users, a static
      header for administration, a per-device bearer secret for hardware. Passwords
      and device secrets are stored as hashes; the gate decision is always
      server-side. \\
    Safety & The gate fails closed. An unreachable backend, a malformed answer or
      an actuator fault never opens the gate; a latched fault requires a reset. \\
    Consistency & State changes are transactional. Session creation is serialised
      on the user row, payment on the session row. \\
    Reliability & Health-checked dependencies in Compose; the test runner waits for
      PostgreSQL and the backend before starting. Stale sessions are reaped so a
      crash cannot block an account indefinitely. \\
    Maintainability & Layered packages (controller, service, repository), DTOs at
      the boundary, constructor injection, no business logic in controllers. \\
    Testability & Every replaceable dependency is an interface
      (\code{PaymentService}, \code{GateControllerClient}, \code{AuthService},
      \code{GateActuator}). 118 automated tests and a coverage gate. \\
    Portability & Backend and tests run in containers; firmware is one Arduino
      sketch with a single configuration header. \\
    Usability & No facility code or location entry: the gate is chosen from a scan.
      Refusals are shown as the backend's own message. \\
    Observability & Public health endpoint, structured gate log lines, JUnit XML
      and JaCoCo reports produced by every run. \\
    Performance & Gate validation is a database read plus a state write; the
      application allows 15\,s for a full BLE round trip including the servo sweep
      and the 8\,s hold. \\
    \bottomrule
  \end{tabularx}
  \caption{Non-functional requirements.}
  \label{tab:nfr}
\end{table}

\section{Software Requirements Specification Summary}

The full specification is maintained in the repository
(\code{dwara-architecture-spec.md}); Table~\ref{tab:srs} is the summary that binds
this report.

\begin{table}[htbp]
  \centering
  \small
  \begin{tabularx}{\textwidth}{@{}lX@{}}
    \toprule
    \textbf{Item} & \textbf{Specification} \\
    \midrule
    Purpose & Phone-controlled parking gate with backend-authorised entry and exit
      and per-second billing. \\
    Product scope & Three deliverables in one repository: Spring Boot backend,
      Expo/React Native application, ESP32 Arduino firmware. \\
    Actors & Driver, attendant/operator, gate controller, scheduled sweep. \\
    External interfaces & REST/JSON over HTTP for the application and the
      administrator; REST/JSON over Wi-Fi for the gate; BLE GATT for the phone; JDBC
      to PostgreSQL. \\
    Data & Users, vehicles, facilities with JSON pricing, sessions, payments, gate
      controllers and commands. All keys are UUIDs; all money is decimal. \\
    Assumptions & One gate per facility; the driver's phone is BLE-capable and runs
      the application in the foreground; the facility has Wi-Fi for the gate. \\
    Constraints & Backend clock is the only billing clock; fixed-timer gate close;
      plaintext BLE in version one; mock payment provider. \\
    Acceptance & Every functional requirement passes its listed evidence; the
      automated suite is green with a 100\% service-layer line-coverage gate; the
      mobile build type-checks; the firmware compiles for ESP32. \\
    Out of scope & Real payments, refunds, occupancy sensing, multiple gates, roles
      beyond \code{USER}, OTA updates, background BLE. \\
    \bottomrule
  \end{tabularx}
  \caption{Software requirements specification summary.}
  \label{tab:srs}
\end{table}

\subsection*{Requirement prioritisation}

All requirements in Table~\ref{tab:functional} were delivered. Had the schedule
slipped, the release order would have been: identity (FR1--FR3), session creation
and the anti-abuse rule (FR7--FR8), gate verdicts and payment (FR10, FR15--FR16),
then history and remote commands. The payment flow and the exit verdict are
inseparable, because a gate that cannot check payment cannot be shipped; history and
administration are the only genuinely optional surfaces.

\section{Hardware and Software Requirements}

\begin{table}[htbp]
  \centering
  \small
  \begin{tabularx}{\textwidth}{@{}lX@{}}
    \toprule
    \textbf{Category} & \textbf{Requirement} \\
    \midrule
    Development host & Docker Engine 24 or later with Compose v2 and \code{bash}.
      No host JDK or Maven is required for the container path (JDK 21 and the Maven
      wrapper are needed only for host-side builds). \\
    Backend runtime & Java 21, Spring Boot 4.1, Spring Security, Spring Data JPA,
      JJWT, Actuator; PostgreSQL 16. \\
    Mobile toolchain & Node.js with npm, Expo 57, React Native 0.86, TypeScript,
      \code{react-native-ble-plx}; a development build on a physical phone for BLE. \\
    Firmware toolchain & \code{arduino-cli} with the ESP32 core, ArduinoJson 7 and
      ESP32Servo; Huge APP partition scheme for the Wi-Fi and BLE build. \\
    Demo tooling & Python 3 with a virtual environment for the serial bridge;
      Postman or Newman for the API contract collection. \\
    Network & Phone and backend on the same LAN, backend port reachable; the gate
      needs Wi-Fi for production validation. \\
    \bottomrule
  \end{tabularx}
  \caption{Software requirements.}
  \label{tab:software-req}
\end{table}

\begin{table}[htbp]
  \centering
  \small
  \begin{tabularx}{\textwidth}{@{}llX@{}}
    \toprule
    \textbf{Part} & \textbf{Quantity} & \textbf{Purpose / note} \\
    \midrule
    ESP32 20-pin development board & 1 & BLE peripheral and Wi-Fi station; 3.3\,V
      logic. \\
    MG90S micro servo & 1 & Gate arm actuator; metal gears, 4.8--6\,V. \\
    External 5\,V supply, at least 1\,A & 1 & Servo power; never the ESP32 3.3\,V
      rail. \\
    Electrolytic capacitor, 470--1000\,\(\mu\)F & 1 & Across the servo supply, close
      to the servo, to absorb stall transients. \\
    Gate arm and hinge & 1 & Light arm; push-through override recommended. \\
    Jumper wires, breadboard, USB cable & --- & Wiring, flashing and serial
      monitor. \\
    \bottomrule
  \end{tabularx}
  \caption{Gate hardware bill of materials.}
  \label{tab:bom}
\end{table}
